\documentclass[12pt]{article}\usepackage[]{graphicx}\usepackage[dvipsnames]{xcolor}
% maxwidth is the original width if it is less than linewidth
% otherwise use linewidth (to make sure the graphics do not exceed the margin)
\makeatletter
\def\maxwidth{ %
  \ifdim\Gin@nat@width>\linewidth
    \linewidth
  \else
    \Gin@nat@width
  \fi
}
\makeatother

\definecolor{fgcolor}{rgb}{0.345, 0.345, 0.345}
\newcommand{\hlnum}[1]{\textcolor[rgb]{0.686,0.059,0.569}{#1}}%
\newcommand{\hlstr}[1]{\textcolor[rgb]{0.192,0.494,0.8}{#1}}%
\newcommand{\hlcom}[1]{\textcolor[rgb]{0.678,0.584,0.686}{\textit{#1}}}%
\newcommand{\hlopt}[1]{\textcolor[rgb]{0,0,0}{#1}}%
\newcommand{\hlstd}[1]{\textcolor[rgb]{0.345,0.345,0.345}{#1}}%
\newcommand{\hlkwa}[1]{\textcolor[rgb]{0.161,0.373,0.58}{\textbf{#1}}}%
\newcommand{\hlkwb}[1]{\textcolor[rgb]{0.69,0.353,0.396}{#1}}%
\newcommand{\hlkwc}[1]{\textcolor[rgb]{0.333,0.667,0.333}{#1}}%
\newcommand{\hlkwd}[1]{\textcolor[rgb]{0.737,0.353,0.396}{\textbf{#1}}}%
\let\hlipl\hlkwb

\usepackage{framed}
\makeatletter
\newenvironment{kframe}{%
 \def\at@end@of@kframe{}%
 \ifinner\ifhmode%
  \def\at@end@of@kframe{\end{minipage}}%
  \begin{minipage}{\columnwidth}%
 \fi\fi%
 \def\FrameCommand##1{\hskip\@totalleftmargin \hskip-\fboxsep
 \colorbox{shadecolor}{##1}\hskip-\fboxsep
     % There is no \\@totalrightmargin, so:
     \hskip-\linewidth \hskip-\@totalleftmargin \hskip\columnwidth}%
 \MakeFramed {\advance\hsize-\width
   \@totalleftmargin\z@ \linewidth\hsize
   \@setminipage}}%
 {\par\unskip\endMakeFramed%
 \at@end@of@kframe}
\makeatother

\definecolor{shadecolor}{rgb}{.97, .97, .97}
\definecolor{messagecolor}{rgb}{0, 0, 0}
\definecolor{warningcolor}{rgb}{1, 0, 1}
\definecolor{errorcolor}{rgb}{1, 0, 0}
\newenvironment{knitrout}{}{} % an empty environment to be redefined in TeX

\usepackage{alltt}
\usepackage{amsmath,amsfonts,amssymb,graphicx,authblk}
\usepackage[font={footnotesize,singlespacing},labelfont=bf]{caption}
\usepackage{titlesec,blkarray, bm} 
\usepackage{float,afterpage}
\usepackage[running,mathlines]{lineno}
\usepackage[vmargin=1in,hmargin=1in]{geometry}
\usepackage[authoryear,sort]{natbib}
\usepackage[dvipsnames]{xcolor}
\usepackage[nodisplayskipstretch]{setspace} 
\usepackage{hyperref}
\usepackage[section]{placeins}
\usepackage{gensymb}
\usepackage{enumitem}
\usepackage{orcidlink}
%\usepackage{hyperref}
\setlist{topsep=.125em,itemsep=-0.15em,leftmargin=0.75cm}
\setlength{\parindent}{0.35in}

\usepackage[sc]{mathpazo} %Like Palatino with extensive math support
\usepackage[subtle]{savetrees}

\usepackage{lineno}
%\renewcommand{\refname}{Literature Cited}
%\renewcommand{\floatpagefraction}{0.9}
%\renewcommand{\topfraction}{0.99}
%\renewcommand{\textfraction}{0.05}

\clubpenalty = 10000
\widowpenalty = 10000

\sloppy 

\usepackage{ifpdf}
\ifpdf
\DeclareGraphicsExtensions{.pdf,.png,.jpg}
\usepackage{epstopdf}
\else
\DeclareGraphicsExtensions{.eps}
\fi

\graphicspath{{/Users/jacobmoutouama/Dropbox/Miller Lab/github/endo-range-limits/Figure}}
%\graphicspath{{C:/Users/tm9/Dropbox/github/ELVI-endophyte-density/Figure}}
\newcommand{\tom}[2]{{\color{red}{#1}}\footnote{\textit{\color{red}{#2}}}}
\newcommand{\jacob}[2]{{\color{blue}{#1}}\footnote{\textit{\color{blue}{#2}}}}
\renewcommand\Authfont{\large}    % sets author names to small
\renewcommand\Affilfont{\footnotesize} % sets affiliations even smaller
%\doublespacing



% I cannot get \orcidlink working on my computer so for now replacing this with \textit

%-------------------------------------------------------------------------
\title{\Large When Mutualism Becomes Costly: Abiotic Stress Shapes Grass Demography}
\author[1]{Jacob K. Moutouama\,\orcidlink{0000-0003-1599-1671} \thanks{Corresponding author: jmoutouama@gmail.com}}
\author[1]{Julia Martin}
\author[1]{Ulisses Rizo}
\author[1]{Malcolm Sherwood}
\author[1]{Emily Chong}
\author[1]{Dajanae Pearson}
\author[1]{Alexandra Jimenez Martín}
\author[2]{Josh Fowler}
\author[1]{Ali Campbell}
\author[1]{Chris Oxley}
\author[1]{Karl Schrader}
\author[1]{Tom E.X. Miller\,\orcidlink{0000-0003-3208-6067}}
\affil[1]{Program in Ecology and Evolutionary Biology, Department of BioSciences, Rice University, Houston, TX USA}
\affil[2]{University of Miami, Department of Biology, Miami, FL USA}
\date{} % clear date
%\renewcommand\Authands{ and }

\sloppy

%-------------------------------------------------------------------------
\IfFileExists{upquote.sty}{\usepackage{upquote}}{}
\frenchspacing
\begin{document}
%\SweaveOpts{concordance=TRUE}
\renewcommand{\baselinestretch}{1.2}
\maketitle
% \bigskip 
%45 character limit on running head
\noindent\textbf{Running header:} Stress and Grass–Endophyte Interactions
\bigskip 

\noindent\textbf{Keywords:} demography, mutualism, parasitism, plant–fungi interactions, range limits, stress gradient hypothesis

\bigskip 
\noindent\textbf{Submitted to:} \textit{Ecology letters} (Letter)

\bigskip 
\noindent\textbf{Data accessibility statement:} 
If the paper is accepted, all data and computer scripts supporting the results will be archived in a Zenodo package, with the DOI included at the end of the article. 
During peer review, our code (Stan and R) is available at \url{https://github.com/jmoutouama/endo-range-limits}. 

\bigskip 
\noindent\textbf{Conflict of interest statement:} None.

\bigskip
\noindent\textbf{Authorship statement:}
%J.K.M., A.C. and T.E.X.M. designed the study.
%A.C. and T.E.X.M. collected the data. 
%All authors conducted the statistical analyses and modeling.
%J.K.M. drafted the manuscript, and T.E.X.M contributed to revisions.

\bigskip
\noindent\textbf{Abstract:}\\
\noindent\textbf{Main Text:}\\
\noindent\textbf{Figures: 6}\\
\noindent\textbf{Tables: X}\\
\noindent\textbf{References: X}

\newpage
%\SweaveOpts{concordance=TRUE}
\linenumbers
%-------------------------------------------------------------------------
\spacing{1.25}
\section*{Abstract}
%150 words limit for Ecology letter. The number now is 
Interactions between plants and fungi are ubiquitous and can strongly influence plant fitness.
Yet, how endophytes affect host demography across abiotic stress gradients and under herbivory remains unclear.
We studied three cool-season grasses hosting seed-transmitted fungal endophytes using geographically distributed field experiments along an aridity gradient, with factorial manipulations of endophyte presence and herbivory. 
Bayesian models showed that all vital rates declined toward more arid sites, indicating reduced performance under stress. 
Declines were more pronounced in individuals hosting endophytes, suggesting that interactions typically beneficial under benign conditions can become costly under abiotic stress. 
Herbivory had no detectable effect on demography. 
Considering both abiotic and biotic factors, abiotic stress was the primary driver of vital rates, with herbivory and endophytes having little effect. 
These results highlight that mutualistic fungi can become liabilities under stress, and that abiotic conditions largely determine species demography across environmental gradients.
%--------------------------------------------------------------------
\newpage
\section*{Introduction}
Plant-microbe symbioses are widespread and ecologically important. 
These interactions are famously context-dependent, where the direction and strength of the interaction outcome depends on the environment in which it occurs \citep{hoeksema2015context, bronstein1994conditional}.
For example, plant-microbe  interactions that are beneficial under stressful conditions  may become neutral or parasitic under benign conditions \citep{giauque2019endophyte}.
Under biotic stress (e.g., herbivory), endophyte symbiosis can benefit host plants by facilitating the production of secondary compounds that deter feeding or exert direct toxicity, thereby reducing insect growth, survival, and oviposition \citep{atala2022fungal,bastias2017epichloe,vega2008insect}.
Similarly under abiotic stress (e.g., low precipitation or high temperature), symbionts can increase their host tolerance \citep{clay2002evolutionary}.
Mycorrhizal fungi form extensive networks that increase root surface area, helping plants absorb more water and nutrients such as phosphorus and nitrogen \citep{amijee1989development}. 
Symbiotic microbes also stimulate plants to produce antioxidant enzymes like superoxide dismutase and catalase that reduce oxidative damage caused by stresses such as drought, salinity, and heat \citep{ruiz1996superoxide}.
 All these mechanisms could help the host species better buffer against environmental stochasticity \citep{fowler2023geographic}.
However, in resource-limited environments, maintaining symbionts can also impose significant costs on hosts \citep{johnson1997functioning, bronstein2001exploitation}. 
Because microbes depend on host-derived carbon and nutrients, plants must allocate part of their limited resources to sustaining the symbiont.
This investment can reduce host growth, reproduction, or competitive ability, particularly when the benefits provided by the symbiont do not outweigh these costs \citep{klironomos2003variation}. 

%\tom{However, in many plant-microbe interactions, host protection is not guaranteed solely by the presence of a symbionts; rather, the density of the symbiont can determine the effectiveness of this protection \citep{laughton2014combined}. 
%Higher endophyte densities may lead to increased resource exploitation by the symbiont, potentially imposing costs on the host, such as reduced growth or reproduction \citep{faeth2009asexual}.}{I like these sentences and the idea is really interesting and important -- but if we are not using the endo density data in this paper then I would remove.}
%This cost is often manifested by a reduction in host biomass or reproductive success \citep{faeth2009asexual}, particularly under harsh abiotic conditions \citep{cui2024review}. 
%Ultimately, these context-dependent costs and benefits may underlie the observed distribution of host species.

Context dependence raises the hypothesis that plant–microbe interactions are likely to vary across abiotic stress gradients, such as high temperatures or low precipitation, as well as under biotic stressors like intense herbivory, with significant implications for host fitness. 
If microbial symbiosis provides greater benefits under these stressful conditions, symbionts could enhance host survival, growth, and reproduction \citep{allsup2023shifting,rudgers2020climate}. 
For example, fungal endophytes enhance host tolerance to abiotic stress in \emph{Bromus laevipes} by improving population survival in dry environments \citep{david2019soil,afkhami2014mutualist}.
Even when symbionts do not directly increase survival under stress, they may still promote host fitness by enhancing growth and reproduction, thereby offsetting the negative effects of lower survival rates \citep{yule2013costs}. 
Conversely, if microbial symbiosis becomes costly under stressful conditions, symbionts could reduce host fitness \citep{benning2021microbes,benning2021plant,bennett2022costs}. 
For instance, mycorrhizal fungi require substantial carbon subsidies from the host plant (up to 20\% of photosynthetically fixed carbon) \citep{thirkell2020carbon,soudzilovskaia2015quantitative}. 
When photosynthesis is limited by abiotic stress, this carbon cost may outweigh the benefits of symbiosis, leading to reduced host fitness.

%Mutualist limitation reduces population fitness and therefore limits range expansion in \emph{Medicago polymorpha} populations \citep{lopez2021microbial}. 
%Although context dependence, along with spatiotemporal variations in abiotic environmental conditions may reduce the effectiveness of the benefits provided by the symbiont to the host species, our understanding of the mechanisms that alter host-symbiont interactions across species geographic range remains limited.

%\tom{Despite growing interest in the ecological and evolutionary roles of symbiosis, most studies focus on the host plant’s perspective and rarely manipulate symbiont prevalence to evaluate how environmental variation shapes the symbiont’s effects on host demography across geographic ranges. 
%Yet, symbionts promote their own fitness by influencing host life history traits and enhancing resistance to stress \citep{kazenel2015mutualistic, giauque2019endophyte, saikkonen1998fungal}. 
%Theory predicts that exclusive vertical transmission fosters mutualism and leads to high symbiont prevalence in host populations \citep{fine1975vectors}. 
%Field studies suggest that, under stressful conditions, beneficial symbionts may increase in prevalence and even reach fixation if their fitness advantages outweigh reproductive costs \citep{donald2021context}. 
%However, populations with lower prevalence may experience higher symbiont loss among offspring \citep{afkhami2008symbiosis}. 
%Thus, overlooking symbiont dynamics could result in inaccurate predictions of host population responses to global change.}{Again, I really like this paragraph, but we need to think about whether and how the symbiont density data will be used here. Is that what you are pointing towards here? Otherwise I am not sure why you are bringing up symbiont prevalence, since I did not think the prevalence data were going into this paper.}

Despite growing interest in the role of symbiosis in host performance, few studies have examined under natural conditions how abiotic and biotic stressors shape the effects of symbionts on host demography along abiotic stress gradients \citep{louthan2015and}. 
Moreover, it remains unclear which stressor is most critical in eliciting the benefits of endophyte symbiosis, or how multiple stressors may interact to influence host demography. 
Herbivory is a biotic factor that can significantly affect host fitness \citep{benning2019biotic} and mediate the outcomes of symbiotic interactions. 
Herbivores reduce host fitness through tissue damage, defoliation, and altered resource allocation \citep{maron2006herbivory}, but they may also shape the selective landscape in ways that favor symbiont persistence or transmission \citep{clay2005herbivores}. 
For example, herbivory can enhance the fitness benefits of endophyte infection by triggering host defense pathways or promoting vertical transmission of the symbiont \citep{gundel2020simulated,agrawal1999transgenerational}. 
If the ecological context of herbivory alters the net fitness effects of symbiosis, then interactions between herbivory and abiotic stress (e.g., drought or high temperature) may play a critical role in determining host population persistence.


Working across an aridity gradient in the south-central US, we investigated how endophyte symbiosis influences host demography under varying abiotic stress. 
This region provides an ideal system to study stressful gradients because it spans a strong east–west precipitation gradient, with mesic sites in the east transitioning to increasingly arid conditions toward the west. 
We also incorporated herbivory as a biotic stressor to evaluate its potential interactive effects with abiotic conditions. 
Specifically, we studied the symbiotic association between three native cool-season grass species (\emph{Agrostis hyemalis}, \emph{Elymus virginicus}, and \emph{Poa autumnalis}) and their vertically transmitted fungal symbionts (\emph{Epichloë} spp.). 
Our experiment was designed to address the following questions:
\begin{enumerate}
    \item Under combined abiotic stress and herbivory, do endophytes provide fitness benefits or impose costs on their host plants?
    \item Between abiotic and biotic factors, which plays a stronger role in shaping host population demography in association with endophytes?
\end{enumerate}



%\begin{enumerate}
%    \item Because endophytes often produce defensive compounds (e.g., alkaloids), plants with endophytes will experience less herbivore damage and maintain higher vital rate compared to plants without endophytes  \tom{when grazed}{We do not have much direct information about `grazing'. I would keep the focus on herbivore access or exclusion, since that is what we directly measured.}.
%    \item Because endophytes can improve \tom{drought tolerance}{I would lead with the abiotic stressors, since that is more central to the experimental design.} (e.g., by enhancing water use efficiency or root growth), plants with endophytes will outperform plants without endophytes more in dry conditions than in wet conditions.
%     \item The combined stresses of herbivory and low precipitation will create conditions where the mutualistic benefits of endophytes are most pronounced, giving plants with endophytes a strong performance advantage over plants without endophytes.
%     \item The combined stresses of herbivory and low precipitation will amplify the physiological demands on the host plant, \tom{such that the additional resource costs of maintaining endophyte associations will  reduce overall plant performance resulting in endophyte-infected plants performing worse than uninfected plants under extreme stress conditions}{I am confused here because this prediction seems to contradict the preceding prediction. But either way, I still think it would be more compelling to phrase these as questions.}.
%\end{enumerate}
%--------------------------------------------------------------------
\section*{Materials and methods}
\subsection*{Study system}
We conducted our study across seven sites spanning a strong east–west precipitation gradient from Louisiana to Texas, USA (Fig.~\ref{fig:site}).  
Louisiana represents mesic conditions, whereas Texas experiences increasingly arid conditions, making this gradient an ideal natural laboratory to study the effects of precipitation and aridity on plant–microbe interactions.  
Mean annual temperature did not differ significantly across sites (Fig.~\ref{Sup:temp}), which is why we focus on precipitation as the primary abiotic stressor.  
This gradient also overlaps the natural distributions of our focal species, ensuring realistic exposure to climatic and biotic conditions (Fig.~\ref{fig:site}).  

At each site, we established eight 1.5 m $\times$ 1.5 m plots, each containing 18 individuals of one of three cool-season perennial grasses: \emph{Agrostis hyemalis}, \emph{Elymus virginicus}, and \emph{Poa autumnalis} \citep{shaw2011guide}.  
These species typically host seed-transmitted fungal endophytes.  
Seeds were collected in summer–fall 2021 from naturally symbiotic populations and either heat-treated to remove symbionts ($E^-$) or left untreated ($E^+$), maintaining the same genetic lineages as the source populations.  

Plots were randomly assigned one of four initial endophyte frequency treatments (80\%, 60\%, 40\%, or 20\%) and one of two herbivory treatments (herbivore access or exclusion).  
%This design was informed by natural patterns of endophyte prevalence in Texas \citep{donald2021context, sneck2017variation, rudgers2009fungus}, and fungal genotyping confirmed pathways for secondary metabolites that may affect host resistance to drought and herbivory \citep{beaudry1951seed}.  
Full details of the study design are provided in the supplementary materials (Supplementary Method \ref{sec:study_design}).

\begin{figure}[H]
\centering
\includegraphics[width=0.85\textwidth]{clim_map.pdf}
\caption{Distribution of common garden sites along the longitudinal aridity gradient in the central and southern Great Plains for: 
A) \emph{Agrostis hyemalis} (AGHY), 
B) \emph{Elymus virginicus} (ELVI), and 
C) \emph{Poa autumnalis} (POAU). 
The first three panels (A–C) are based on 30-year normals of climate predictions (1993–2023) from PRISM data, which illustrate the precipitation gradient across sites (scale on the right is in mm). 
Red dots indicate the locations of source populations, while grey dots represent Global Biodiversity Information Facility (GBIF) occurrence records for each species within the study area. 
D) Cumulative precipitation measured during the study period at each site, confirming the precipitation gradient shown in panels A–C.}
\label{fig:site}
\end{figure}

\subsection*{Demographic data collection}
We collected demographic data including survival, growth, and reproduction in May and June of 2023, 2024 and 2025, which coincided with the flowering season of \emph{Agrostis hyemalis}, \emph{Elymus virginicus}, and \emph{Poa autumnalis}. 
For each individual, plant survival was recorded as a binary variable (alive or dead), and plant size was measured as the number of living tillers.  
We recorded the number of inflorescences per plant and counted the number of spikelets on up to three inflorescences from three reproducing plants.  
Spikelet counts were limited to three reproducing tillers per plot due to the time-consuming nature of this measurement.  
We used the number of spikelets from these three tillers to estimate the average number of spikelets per plant.

\subsection*{Modeling the effects of aridity and herbivory on endophyte–host demography}
\label{sec:aridity_herbivory}
To quantify how abiotic stress and herbivory shape plant–fungal interactions, we developed hierarchical models for each vital rate (survival, growth, and fertility). 
Survival was modeled with a Bernoulli distribution, growth with a Gaussian distribution, and fertility, measured as the number of inflorescences and the number of spikelets, with negative binomial distributions.
All species were modeled jointly to estimate pooled fixed effects. 
Each vital rate model included normally distributed random effects to account for site-to-site variation ($\phi_{\text{site}[i]} \sim \mathcal{N}(0,\; \sigma_{\text{site}})$), source-to-source variation ($\omega_{\text{source}[i]} \sim \mathcal{N}(0,\; \sigma_{\text{source}})$) associated with the provenance of transplants used in the common garden, and plot-to-plot variation ($\rho_{\text{plot}[i]} \sim \mathcal{N}(0,\; \sigma_{\text{plot}})$).
By modeling species jointly with this hierarchical structure, we increase statistical power to detect shared effects while still capturing species-specific responses.
To allow for non-linear effects of precipitation, we included a quadratic term, reflecting the expectation that vital rates may peak at intermediate precipitation and decline at both low and high levels. 
The expected value of each vital rate for individual $i$ was modeled as:
\begin{equation}\label{eq:candidates}
\begin{aligned}
\mu_i &= \beta_0 + \beta_P P_i + \beta_{endo}\,endo_i + \beta_{herb}\,herb_i \\
&\quad + \beta_{endo \times P}\,endo_i P_i 
+ \beta_{herb \times endo}\,herb_i endo_i 
+ \beta_{herb \times P}\,herb_i P_i \\
&\quad + \beta_{P^2} P_i^2 + \beta_{endo \times P^2}\,endo_i P_i^2 \\
&\quad + \phi_{\text{site}[i]} + \omega_{\text{source}[i]} + \rho_{\text{plot}[i]} 
\end{aligned}
\end{equation}
\noindent where $P$ is precipitation, $\text{endo}$ indicates endophyte presence, and $\text{herb}$ indicates herbivory. 
The intercept $\beta_0$ represents the grand mean vital rate across species. 
The coefficients $\beta_P$ and $\beta_{P^2}$ describe linear and quadratic responses to precipitation, respectively, while $\beta_{\text{endo}}$ and $\beta_{\text{herb}}$ capture the main effects of endophytes and herbivory. 
The interaction terms capture combined effects: $\beta_{\text{endo}\times P}$ captures how endophytes alter the response to precipitation, $\beta_{\text{herb}\times P}$ captures how herbivory affects the response to precipitation, $\beta_{\text{herb}\times \text{endo}}$ represents the joint effect of herbivory and endophytes, and $\beta_{\text{endo}\times P^2}$ captures how endophytes influence the quadratic response to precipitation.
%Models were implemented in R \citep{RCoreTeam} using Stan via the RStan interface \citep{Rstan}.

\subsection*{Disentangling abiotic and biotic drivers of endophyte–host demography}
To disentangle the relative contributions of abiotic and biotic factors to endophyte–host demography, we developed three candidate models. 
The first included only abiotic predictors (precipitation), the second included only biotic predictors (herbivory and endophyte status), and the third included both abiotic and biotic predictors. 
All models followed the same hierarchical structure described above, differing only in the inclusion of predictor terms (Supplementary Method \ref{sec:models}).

To determine the best-supported model for each vital rate, we compared candidate models using leave-one-out cross-validation (LOOCV) \citep{vehtari2017practical}. 
LOOCV combines validation and training by iteratively holding out one observation as a test set while fitting the model on the remaining $n-1$ observations. 
This process is repeated $n$ times, once for each observation, resulting in $n$ fitted models \citep{silva2024robust}. 
The overall test error estimate is then obtained by averaging the prediction errors across all $n$ models (Eq.~\ref{eq:Loo}).

\begin{align}\label{eq:Loo}
CV_{n} = \frac{1}{n} \sum_{i=1}^{n} (y_{i} - \hat{y}_{i})^2
\end{align}

All models were implemented in R \citep{RCoreTeam} using Stan via the \texttt{rstan} interface \citep{Rstan}.


%%%%%%%%%%%%%%%%%%%%%
% Acknowledgments
%%%%%%%%%%%%%%%%%%%%%
% You may wish to remove the Acknowledgments section while your paper 
% is under review as the Acknowledgments may reveal your identity.
% If you remove this section, you will need to add it back in to your
% final files after acceptance.
\section*{Results}
\subsection*{Endophyte presence influenced host survival and reproduction across the abiotic stress gradient, with no effects of herbivory}
We used Bayesian mixed-effect models to test the effects of abiotic stress (low vs. high precipitation), endophyte association (E$^{-}$ vs. E$^{+}$), and herbivore exclusion (fenced vs. unfenced) on survival, growth, and reproductive traits in three cool-season grass species: \textit{Agrostis hyemalis}, \textit{Elymus virginicus}, and \textit{Poa autumnalis}. All models effectively captured key aspects of the data, including means, standard deviations, and quantiles (Fig.~\ref{sup:ppc_vital}).  
High abiotic stress reduced survival, with E$^{-}$ individuals outperforming E$^{+}$ individuals, indicating a cost of endophyte association under drought conditions (Fig.~\ref{fig:Surv}). Growth was generally similar across species and treatments, except in \textit{Agrostis hyemalis}, where E$^{+}$ plants showed higher relative growth along the stress gradient (Fig.~\ref{fig:growth}). Inflorescence production differed little between E$^{+}$ and E$^{-}$ individuals, except for \textit{Poa autumnalis}, which produced more inflorescences under low abiotic stress (high precipitation; Fig.~\ref{fig:flowering}). Spikelet production also reflected costs of endophyte association under high stress, particularly in \textit{Elymus virginicus}, where E$^{+}$ individuals produced fewer spikelets than E$^{-}$ individuals (Fig.~\ref{fig:spikelet}).  
Herbivore exclusion had minimal effects on all vital rates, with slightly higher survival and spikelet production in fenced plots, indicating that precipitation and endophyte status were the primary drivers of variation in host demography.

\begin{figure}[H]
\centering
\includegraphics[width=0.8\textwidth]{PrSurvival_diff.pdf}
\caption{XXX 
}
\label{fig:Surv}
\end{figure}

\begin{figure}[H]
\centering
\includegraphics[width=0.8\textwidth]{Growth_diff.pdf}
\caption{Predicted relative growth (top panels) and endophyte effect (\(\Delta\) E+ - E-, bottom panels) across precipitation gradients. Lines show posterior means, shading shows 90\% credible intervals, and points indicate mean observed growth per plot. \textbf{Colors}: E+ (blue), E- (red). \(\Delta = 0\) (dashed line) indicates no endophyte effect. Panels are organized by species (AGHY, ELVI, POAU) and herbivory treatment (Fenced, Unfenced).}
\label{fig:growth}
\end{figure}



\begin{figure}[H]
\centering
\includegraphics[width=0.8\textwidth]{Flowering_diff.pdf}
\caption{XXX 
}
\label{fig:flowering}
\end{figure}

\begin{figure}[H]
\centering
\includegraphics[width=0.8\textwidth]{Spik_ppt_diff.pdf}
\caption{XX
}
\label{fig:spikelet}
\end{figure}



\section*{Acknowledgements}
This research was supported by National Science Foundation Division of Environmental Biology awards.



\newpage
\bibliographystyle{apalike}
\bibliography{References}
%--------------------------------------------------------------------
\newpage
\clearpage 
\setcounter{equation}{0}
\setcounter{figure}{0}
\setcounter{section}{0}
\setcounter{table}{0}
\renewcommand{\theequation}{S.\arabic{equation}}
\renewcommand{\thetable}{S-\arabic{table}}
\renewcommand{\thefigure}{S-\arabic{figure}}
\renewcommand{\thesection}{S.\arabic{section}}

\centerline{\Large{\textbf{Supporting Information}}}
\section{Supplementary Methods} 
\subsection {Study design}\label{sec:study_design}
\subsubsection {Abiotic stress}
Common gardens were established at seven sites spanning the geographic range of \emph{Elymus virginicus} from Louisiana to Texas, USA (Fig.~\ref{sup:md}). Sites were selected along a pronounced aridity gradient, with Louisiana representing mesic conditions and Texas increasingly arid conditions. Mean annual temperature did not differ significantly among sites (Fig.~\ref{Sup:temp}), confirming that precipitation is the primary abiotic gradient. Sites were chosen to reflect natural environmental conditions, including shaded areas under tree canopy or near shrubs.

\subsubsection {Plot establishment}
At each site, we established eight plots measuring 1.5 m $\times$ 1.5 m. Existing vegetation was removed by tilling to reduce competition. Within each plot, 15 individuals of \emph{Elymus virginicus} were planted approximately 15 cm deep in an evenly spaced 4 $\times$ 4 grid, with positions randomly assigned to reduce positional effects.

\subsubsection{Endophyte and herbivory treatments}
Seeds were collected in summer–fall 2021 from naturally symbiotic populations (source populations). Seeds were either heat-treated to remove endophytes ($E^-$) or left untreated ($E^+$), maintaining the same genetic lineages. Each plot was randomly assigned a starting endophyte frequency (80\%, 60\%, 40\%, or 20\%) and a herbivory treatment. For herbivore exclusion plots, plants were enclosed with 1.2 m tall mesh fencing and sprayed with insecticide to prevent vertebrate and insect herbivory. For herbivore access (control) plots, plants were half-enclosed with mesh netting. All plots received comparable quantities of source population material.

%\subsubsection{Initial measurements and monitoring}
%After planting, all plots were watered to promote establishment. For each individual plant, we recorded initial tiller counts, flowering status, plot position, endophyte status, and source population.
%
\subsubsection{Experimental rationale}
This experimental design was informed by natural patterns of endophyte prevalence in Texas \citep{donald2021context, sneck2017variation, rudgers2009fungus}. Fungal genotyping confirmed the presence of biosynthetic pathways for secondary metabolites, such as peramine, loline, and ergot alkaloids, which may enhance host resistance to drought and herbivory \citep{beaudry1951seed}. The design allows testing how endophyte frequency and herbivory interact with abiotic stress (aridity) to influence host demography from core to edge of the species range.

\subsubsection{Source populations and identification of individual endophyte status}
Plants used in the common garden experiment were derived from natural populations throughout the south-central US. Seeds were collected from these populations, and some were heat-treated at 60°C for approximately five days (120 hours) to produce endophyte-negative plants ($E^-$), a method that eliminates endophytes without affecting seed viability. All seeds, both heat-treated and non-heat-treated, were planted in the Rice University greenhouse, where seedlings were fertilized every two weeks. Seedlings were then vegetatively propagated to produce enough individuals for the experiment (N = 840). Before field planting, we confirmed endophyte status ($E^+$ or $E^-$) of all seedlings using either microscopy or an immunoblot assay. Leaf tissues were stained with aniline blue lactic acid and viewed under a compound microscope at 200x–400x to identify fungal hyphae. The immunoblot assay (Phytoscreen field tiller endophyte detection kit, Agrinostics Ltd. Co.) uses monoclonal antibodies that target proteins of \emph{Epichloë} spp. and a chromagen to visually indicate presence or absence. Both methods yield similar detection rates.

\subsection{Candidate models for abiotic and biotic drivers} \label{sec:models}

To disentangle the relative contributions of abiotic and biotic factors, we developed three candidate models:

\begin{enumerate}
    \item \textbf{Abiotic model:} includes precipitation ($P$) and its quadratic term ($P^2$) as predictors. 
    The expected value of each vital rate for individual $i$ is:
    \begin{align*}
    \mu_i &= \beta_0 + \beta_P P_i + \beta_{P^2} P_i^2 \\
    &\quad + \phi_{\text{site}[i]} + \omega_{\text{source}[i]} + \rho_{\text{plot}[i]}
    \end{align*}

    \item \textbf{Biotic model:} includes endophyte status ($\text{endo}$), herbivory ($\text{herb}$), and their interaction ($\text{endo} \times \text{herb}$):
    \begin{align*}
    \mu_i &= \beta_0 + \beta_{\text{endo}}\,endo_i + \beta_{\text{herb}}\,herb_i 
            + \beta_{\text{endo}\times \text{herb}}\,endo_i\,herb_i \\
    &\quad + \phi_{\text{site}[i]} + \omega_{\text{source}[i]} + \rho_{\text{plot}[i]}
    \end{align*}

    \item \textbf{Full model:} combines both abiotic and biotic predictors, including all relevant interactions:
    \begin{align*}
    \mu_i &= \beta_0 + \beta_P P_i + \beta_{P^2} P_i^2 + \beta_{\text{endo}}\,endo_i + \beta_{\text{herb}}\,herb_i \\
    &\quad + \beta_{\text{endo}\times P}\,endo_i P_i 
          + \beta_{\text{herb}\times P}\,herb_i P_i 
          + \beta_{\text{herb}\times \text{endo}}\,herb_i\,endo_i 
          + \beta_{\text{endo}\times P^2}\,endo_i P_i^2 \\
    &\quad + \phi_{\text{site}[i]} + \omega_{\text{source}[i]} + \rho_{\text{plot}[i]}
    \end{align*}
\end{enumerate}

All models included a hierarchical random effects structure accounting for site, source, and plot variation, with $\phi_{\text{site}[i]} \sim \mathcal{N}(0, \sigma_{\text{site}})$, $\omega_{\text{source}[i]} \sim \mathcal{N}(0, \sigma_{\text{source}})$, and $\rho_{\text{plot}[i]} \sim \mathcal{N}(0, \sigma_{\text{plot}})$.

\section {Supporting Tables}

\begin{table}[H]
\centering
\caption{Populations and Coordinates}
\begin{tabular}{|p{2in}|l|l|l|}
\hline
\textbf{Population} & \textbf{Location} & \textbf{Coordinates} & \textbf{NE+ / NE-} \\ \hline
Sam Houston State University Field Station  & Huntsville, TX & 30°44'30.5"N, 95°28'28.2"W & NE+ = 39, NE- = 50 \\ \hline
Texas Tech University  & Junction, TX & 30°28'18.2"N, 99°47'01.7"W & NE+ = 49, NE- = 83 \\ \hline
Palmetto State Park  & Palmetto, TX & 29.591623, -97.584781 & NE+ = 250, NE- = 192 \\ \hline
Jean Lafitte National Historical Park and Preserve  & Jean Lafitte, LA & 29.785105, -90.114933 & NE+ = 83, NE- = 94 \\ \hline
\end{tabular}
\end{table}

\begin{table}[H]
\caption{Model comparison by vital rate using ELPD (Expected Log Pointwise Density) differences. The values in the \textbf{elpd\_diff} column represent the difference in the model's ELPD relative to the best model for each vital rate, with positive values indicating a worse model fit. The \textbf{se\_diff} column shows the standard error of the ELPD difference. Best models for each vital rate are bolded. Model selection is based on the ELPD difference and associated standard error, with smaller ELPD differences and lower standard errors indicating better model fit.}
\centering
\begin{tabular}{|l|l|r|r|}
\hline
\textbf{Vital rate} & \textbf{Model} & \textbf{elpd\_diff} & \textbf{se\_diff} \\
\hline
Survival   & \textbf{model1} & \textbf{0.0}   & \textbf{0.0} \\
           & model2          & -1.0           & 1.5          \\
           & model3          & -1.0           & 1.5          \\
\hline
Growth     & \textbf{model3} & \textbf{0.0}   & \textbf{0.0} \\
           & model2          & -1.3           & 2.0          \\
           & model1          & -2.1           & 2.2          \\
\hline
Flowering  & \textbf{model2} & \textbf{0.0}   & \textbf{0.0} \\
           & model3          & 0.0            & 3.2          \\
           & model1          & -1.3           & 3.9          \\
%\hline
%Spikelet   & \textbf{model2} & \textbf{0.0}   & \textbf{0.0} \\
%           & model3          & -4.6           & 4.2          \\
%           & model1          & -6.4           & 5.3          \\
\hline
\end{tabular}
\label{tab:elpd_vital_rates}
\end{table}





\section {Supporting Figures}

\begin{figure}[H]
		\centering
		\includegraphics[width=1\linewidth]{Climate_prism_2023_2025.pdf}
		\caption{Abiotic conditions across the seven study sites. 
		(A) Mean annual precipitation showing the pronounced east–west aridity gradient from mesic Louisiana to arid Texas. 
		(B) Mean annual temperature across sites, which does not differ much. 
		Panels highlight that precipitation, rather than temperature, constitutes the primary abiotic gradient in this system.}
		\label{Sup:temp}
\end{figure}
\clearpage  % forces the next figure to start on a new page



\begin{figure}[H]
\centering
\includegraphics[width=1\textwidth]{combined_ppc_plots.pdf}
\caption{
Posterior predictive checks for the full demographic model. 
Dark blue lines show observed data ($y$), and light blue lines show replicated datasets generated from posterior predictions ($y_{\text{rep}}$). 
The strong overlap across survival, growth, flowering, and spikelet production indicates adequate model fit for all vital rates.
}
\label{sup:ppc_vital}
\end{figure}

%
%\begin{figure}[H]
%		\centering
%		\includegraphics[width=0.50\linewidth]{ppt_surv_coeff.pdf}
%		\caption{Coefficient estimates from Bayesian models linking survival and geographic distance.
%		 Points represent posterior means and lines indicate 95\% credible intervals.
%		 Species abbreviations: \textit{Agrostis hyemalis} (AGHY), \textit{Elymus virginicus} (ELVI), and \textit{Poa autumnalis} (PAOU).
%		 Horizontal dashed lines denote zero effect.}
%		\label{Sup:temp_variation}
%\end{figure}
%
%\begin{figure}[H]
%		\centering
%		\includegraphics[width=0.50\linewidth]{ppt_grow_coeff.pdf}
%		\caption{Coefficient estimates from Bayesian models linking growth and geographic distance.
%		 Points represent posterior means and lines indicate 95\% credible intervals.
%		 Species abbreviations: \textit{Agrostis hyemalis} (AGHY), \textit{Elymus virginicus} (ELVI), and \textit{Poa autumnalis} (PAOU).
%		 Horizontal dashed lines denote zero effect.}
%		\label{Sup:temp_variation}
%\end{figure}
%
%\begin{figure}[H]
%		\centering
%		\includegraphics[width=0.50\linewidth]{ppt_flow_coeff.pdf}
%		\caption{Coefficient estimates from Bayesian models linking flowering and geographic distance.
%		 Points represent posterior means and lines indicate 95\% credible intervals.
%		 Species abbreviations: \textit{Agrostis hyemalis} (AGHY), \textit{Elymus virginicus} (ELVI), and \textit{Poa autumnalis} (PAOU).
%		 Horizontal dashed lines denote zero effect.}
%		\label{Sup:temp_variation}
%\end{figure}
%
%\begin{figure}[H]
%		\centering
%		\includegraphics[width=0.50\linewidth]{ppt_spik_coeff.pdf}
%		\caption{Coefficient estimates from Bayesian models linking flowering and geographic distance.
%		 Points represent posterior means and lines indicate 95\% credible intervals.
%		 Species abbreviations: \textit{Agrostis hyemalis} (AGHY), \textit{Elymus virginicus} (ELVI), and \textit{Poa autumnalis} (PAOU).
%		 Horizontal dashed lines denote zero effect.}
%		\label{Sup:temp_variation}
%\end{figure}
%
%

%\begin{figure}[H]
%		\centering
%		\includegraphics[width=0.95\linewidth]{Figures/fig_pr_past_present_future.pdf}
%		\caption{Precipitation variation across the study sites from 1990 to 2100.
%		ds: Dormant season, dg: Growing season.}
%		\label{Sup:pr_variation}
%\end{figure}
%
%\begin{figure}[H]
%		\centering
%		\includegraphics[width=0.95\linewidth]{Figures/MIROC.pdf}
%		\caption{Past, Observed, present and future (MIROC Model) climate data across the study area.}
%		\label{Sup:projectionMIROC}
%\end{figure}
%
%\begin{figure}[H]
%		\centering
%		\includegraphics[width=0.99\linewidth]{Figures/ACCESS.pdf}
%		\caption{Past, Observed, present and future (ACCESS Model) climate data across the study area.}
%		\label{Sup:projectionACCESS}
%\end{figure}
%
%\begin{figure}[H]
%		\centering
%		\includegraphics[width=0.99\linewidth]{Figures/CESM1.pdf}
%		\caption{Past, Observed, present and future (CESM1 Model) climate data across the study area.}
%		\label{Sup:projectionCESM1}
%\end{figure}
%
%\begin{figure}[H]
%		\centering
%		\includegraphics[width=0.99\linewidth]{Figures/CMCC.pdf}
%		\caption{Past, Observed, present and future (CMCC Model) climate data across the study area.}
%		\label{Sup:projectionCMCC}
%\end{figure}
%
%\begin{figure}[H]
%		\centering
%		\includegraphics[width=0.959\linewidth]{Figures/site_year_weather.pdf}
%		\caption{Climate variation across the study sites during the monitoring period.}
%		\label{Sup:climate_variation}
%\end{figure}


\end{document}
